\section{Methodology}\label{sec:pre_confirmation:methodology}
The research programme of the thesis is situated within the constructive research tradition of computer science: its primary contribution is a designed and built artefact rather than an empirical observation or theoretical derivation.
Design Science Research (DSR)~\cite{Hevner2004DesignScienceIS} is the governing methodology.
DSR is the established framework for research whose contribution is a novel IT artefact evaluated against a practical problem.
It provides both a process model for artefact construction and an evaluative standard grounded in utility rather than optimality~\cite{Hevner2004DesignScienceIS}.
The research programme satisfies that condition: the planned contribution is a quantum-implicit programming language and its supporting formal and computational infrastructure, evaluated against the abstraction gap identified in Section~\ref{sec:pre_confirmation:literature_review}.

Three alternative research designs were considered and rejected.
Empirical and experimental methods require a pre-existing system to observe or experiment upon; no such system exists prior to this work.
A proof-only approach would establish the formal properties of the type system and compilation pipeline but cannot account for the design and construction of the language and compiler as research contributions.
Qualitative methods are not applicable to artefact construction.
The most substantive objection to the DSR design is that improved compilation technology renders a new source language unnecessary.
If a compiler can automatically infer quantum-level decisions, a quantum-implicit language provides no additional value.
This objection fails because a compiler can only infer what the source language provides.
Quantum-explicit languages encode quantum-level decisions at the vocabulary level, leaving no higher-level intent for a compiler to recover.
DSR explicitly supports the iterative design-build-evaluate cycle appropriate for language and compiler development and is therefore the appropriate choice.

The research follows a problem-centred entry point into the DSRM process~\cite{Peffers2007DSRM}.
The abstraction gap in quantum programming languages is established as the motivating problem before solution goals are defined.
The six DSRM activities map directly onto the planned research structure.
Problem identification and motivation are addressed through the literature review and the planned systematic review publication, which together establish the abstraction gap and supply the requirements the artefact must satisfy.
The objectives of a solution are defined as a quantum-implicit programming language: a programming abstraction model compilable to gate-based quantum hardware and accessible to programmers without quantum expertise.
Design and development will be delivered through the planned language specification, formal theoretical basis and compilation framework.
Demonstration will be provided through proof-of-concept programs and compiler output on simulators and NISQ hardware.
Evaluation will be conducted against the six criteria defined below.
Communication will be delivered through the four planned incorporated publications and the thesis.
The planned formal proofs will span the design and development and evaluation activities simultaneously.
They will constitute both a component of the design deliverable and the primary evaluation mechanism for SRQ2 and SRQ3.

The planned research contribution comprises three integrated artefact components.
The \textit{language specification} is the primary design contribution, targeting a formal definition of the syntax, semantics and type system of the quantum-implicit language.
The \textit{formal theoretical basis} targets type safety through progress and preservation proofs for the core language.
It also targets semantic preservation for a defined core fragment of the compilation pipeline.
The \textit{compilation framework} targets the lowering of programs in the source language to executable form on gate-based NISQ hardware.
The research adopts a gate-based quantum computing model throughout.
Measurement-based quantum computing, hardware-specific scheduling, pulse-level control, quantum error correction at the language level and full compiler verification are out of scope and identified as future work.

DSR evaluates artefacts against utility for the stated problem~\cite{Hevner2004DesignScienceIS}.
Six criteria operationalise this standard for the planned artefact.
\textbf{Type safety}: the research targets a formal proof of progress and preservation for the core language, with peer review sought as part of the language specification publication.
\textbf{Compilation correctness}: the research targets a formal semantic preservation argument for a defined core fragment, with scope boundaries and limitations stated.
\textbf{Executability}: the research targets compilation and execution of programs on gate-based NISQ hardware via the University of Melbourne's membership of the IBM Quantum Network~\cite{UoMIBMQuantumHub}.
Simulator execution is the minimum threshold.
\textbf{Expressiveness}: the research targets implementation and compilation of at least two non-trivial quantum algorithms without escape hatches.
\textbf{Programming language abstraction model}: the research targets demonstrating that representative algorithmic tasks are expressible using programming language abstractions only, without requiring quantum vocabulary.
\textbf{Absence of mandatory quantum vocabulary}: the research targets case studies and example programs demonstrating that a programmer with a classical background can engage with the language without quantum-mechanical knowledge.
A formal usability study is identified as future work.